% === IX. CONCLUSION ==========================================================
\section{Conclusion and Future Work}
\label{sec:conclusion}

This paper studied the identifiability of a learning-based on-track Pacejka
identification pipeline~\cite{ontrack2025} at full vehicle scale, using a CARLA
simulation architecture that exposes the plant's own per-wheel forces as ground
truth. Transplanted unchanged from a 1:10 platform, the pipeline fails in a
specific, diagnosable and correctable way: a loss of excitation created by its
own virtual-data step, which no change of solver, bounds or tuning removes, and
whose reachable friction depends on the rollout configuration and not the tire.
Two further properties of the loop follow from the same analysis. Stage~3
reconstructs its forces from the rollout's own quasi-steady yaw balance, so its
residual depends on $B$, $C$ and $D$ only through their product and the peak
factor is not recoverable from the rollout at any excitation; it must be
measured on the recorded trace and held. And each fit is handed the trajectory
its own previous answer generated, so stage~4 is a positive feedback path
rather than a contraction and is reformulated as a relaxed fixed-point
iteration at $\beta=\num{0.20}$. The correction ships with an extrapolation
bound on the steering sweep, a frozen regularisation target, identification
against the achieved road-wheel angle, measured mass and geometry, and a
dynamics residual normalised to the units of the data term, and it produces
accepted identifications with no railed coefficient and the correct understeer
sign.

A second surface settles what a single one cannot. With the plant's friction
driven down under the moving vehicle on the linear-decay protocol
of~\cite{llampc2025} until half the grip is gone, the corrected pipeline
recovers the decay rate to \SI{7}{\percent} and the level to
\SIrange{10.9}{17.4}{\percent} once one cycle of transport lag is accounted
for; the falling peak factor re-clamps the reference from \SI{0.51}{\g} of peak
lateral demand to \SI{0.30}{\g}, buying three completed laps at unchanged
tracking error for \SI{18}{\percent} of lap time. The inherited configuration
leaves the track in lap~2, holding a \SI{2.00}{\g} grip ceiling against a plant
offering \SI{0.53}{\g}. Both outcomes have one cause: a peak factor resting on
a box constraint cannot report that anything changed, so the identifier's
failure stays invisible until the controller acts on it.

Two qualifications stand on those numbers: the peak factor is identified by the
brush estimator on the recorded trace, not by the Pacejka fit, which cannot
identify it (Section~\ref{subsec:bcd}); and each configuration is a single run
per surface, so the closed-loop margins are observations rather than estimates
with a confidence interval. We further argue that per-coefficient box
constraints are not a sufficient admissibility criterion for handing a model to
a model-based controller, and place gates on derived physical quantities at
that interface instead.

\subsection{Future work}

\begin{enumerate}
\item \textbf{Shorten the cycle, or step the friction.} The decay runs expose
one cycle of transport lag as the dominant error on a moving surface, worth
\SIrange{10.9}{17.4}{\percent} at the rate used here and proportionally more at
a faster one. Two experiments follow directly, both scenario-file rows: the step
schedules, on which the same lag appears as a bounded overshoot rather than a
standing offset, and a shorter re-identification period, which the
\SI{20.9}{\second} cycle cost of Section~\ref{subsec:runtime} leaves room for.
\item \textbf{Run the ablations of Section~\ref{subsec:ablations} in the loop.}
The relaxation sweep, the $D$-pinning comparison and the S4D and objective
ablations are offline replays, enough to attribute the result but not to confirm
the ordering survives closed-loop coupling. Each is a scenario-file row.
\item \textbf{Excite the peak, and design for it.} The lap reaches
$P_{99}(|\alpha|)=\SI{2.7}{\degree}$ against a plant whose curve peaks at
\SI{8.0}{\degree}, which is where the brush estimator's \SI{11}{\percent} bias
comes from. Equation~\eqref{eq:mureach} inverts into a requirement on the
manoeuvre, so a controller that inserts bounded slalom or ramp-steer excitation
when the identifier reports insufficient slip-angle coverage would move the
pipeline toward active, safety-bounded experiment
design~\cite{goodwin1977design}.
\item \textbf{Hardware replication.} The decisive test of every claim in
Section~\ref{sec:disc} is a physical Formula Student vehicle. The identification
code is unchanged between simulation and hardware; what changes is the
availability of ground truth, which is what makes the simulation study a
necessary precursor rather than a substitute.
\end{enumerate}
